\subsection{Analytical derivation of the spectral ratio on $S^3$}
  \label{sec:analytical_ratio}

  The numerical convergence of $\lambda_2/\lambda_1$ toward $8/3$
  admits an exact analytical counterpart on the unit three-sphere.

  \paragraph{Scalar spectrum of $S^3$.}
    For the positive Laplace--Beltrami operator
    $\Delta_{S^3}=-\mathrm{div}\,\mathrm{grad}$,
    the eigenvalues are~\cite{Berger1971}
    \begin{equation}
      \lambda_k = k(k+2), \qquad k=0,1,2,\ldots,
      \label{eq:S3spectrum}
    \end{equation}
    with multiplicity $(k+1)^2$.
    Hence
    \[
      \lambda_1=3, \qquad \lambda_2=8,
    \]
    and
    \begin{equation}
      \frac{\lambda_2}{\lambda_1}=\frac{8}{3}.
      \label{eq:exact_ratio}
    \end{equation}
    This follows algebraically from the $\mathrm{SU}(2)\cong S^3$
    representation structure and is independent of coordinates,
    discretization, or fibration.

  \paragraph{Role of the Hopf structure.}
    The Hopf fibration $S^1\hookrightarrow S^3\to S^2$
    does not modify the intrinsic spectrum.
    It provides only a geometric interpretation of the
    $k=1$ eigenspace in terms of fiber and base excitations.
    The ratio $8/3$ is therefore an intrinsic spectral
    property of $S^3$, not a consequence of fiber weighting.

  \paragraph{Relation to the numerical estimator.}
    The simulations evaluate a kernel-weighted ratio
    $R(\alpha)=E_{\mathrm{fiber}}(\alpha)/E_{\mathrm{base}}(\alpha)$.
    In the isotropic regime ($\alpha\le 0$), $R(\alpha)$
    reflects geometric partitioning.
    As $\alpha$ increases and fiber modes are selectively excited,
    the normalized fiber response converges to $8/3$.
    This is consistent with~\eqref{eq:exact_ratio}:
    strong fiber excitation isolates the first two distinct
    spectral levels, whose eigenvalue ratio equals $8/3$.

  \paragraph{Diagnostic status.}
    The value $8/3$ is a topological spectral invariant of $S^3$.
    It is not universal for arbitrary relational graphs.
    Its role is diagnostic: substrates reproducing this ratio
    in a projectable regime are spectrally compatible with
    $S^3$-type geometry.
    Different topologies yield different ratios, providing a
    structural discriminator stronger than metric reconstruction alone.
